\begin{examplebox}
    \textbf{\textcolor{CExample}{例一（教材原型：首一四次多项式，试根与长除法）}}
    \textbf{\textcolor{CExample}{题目：}}
    分解 $f(x)=x^4-x^3-3x^2+x+2$，并求方程 $f(x)=0$ 的所有有理根。
    \begin{description}[style=nextline, leftmargin=2.8em, labelsep=0.5em, itemsep=0.52em]
        \item[\textbf{第一步（有理根候选）：}]
            $a_4=1$，$a_0=2$。有理根定理给出 $q\mid1$，故 $q=1$；$p\mid2$，故 $p=\pm1,\pm2$。因此只需依次检验 $\pm1,\pm2$。
            \textbf{理由：} 首项系数为 $1$ 的整系数多项式，有理根只能是整数。
        \item[\textbf{第二步（先试 $x=1$）：}]
            \[
              f(1)=1-1-3+1+2=0.
            \]
            因此 $x-1$ 是因式。
            \textbf{理由：} 代入等于零才能确认候选根成立。
        \item[\textbf{第三步（长除法逐项降次）：}]
            把 $x^4-x^3-3x^2+x+2$ 除以 $x-1$：
            \begin{itemize}[leftmargin=*, itemsep=0.28em, topsep=0.2em]
               \item $x^4\div x=x^3$；减去 $(x-1)x^3=x^4-x^3$，余 $-3x^2+x+2$。
               \item $-3x^2\div x=-3x$；减去 $(x-1)(-3x)=-3x^2+3x$，余 $-2x+2$。
               \item $-2x\div x=-2$；减去 $(x-1)(-2)=-2x+2$，余数为 $0$。
            \end{itemize}
            所以
            \[
              f(x)=(x-1)(x^3+0x^2-3x-2).
            \]
            \textbf{理由：} 每次选一个商式项，使被除式当前最高次项消去；缺少的 $x^2$ 项系数为零。
        \item[\textbf{第四步（对三次商式继续试根）：}]
            令 $g(x)=x^3-3x-2$，计算 $g(-1)=-1+3-2=0$，故 $x+1$ 是 $g(x)$ 的因式。再除一次得
            \[
               g(x)=(x+1)(x^2-x-2)=(x+1)^2(x-2).
            \]
            \textbf{理由：} 一次试根只保证找到一个因式，降次后还要继续求根。
        \item[\textbf{第五步（汇总并检查）：}]
            \[
              \boxed{f(x)=(x-1)(x-2)(x+1)^2.}
            \]
            根为 $1,2,-1$，其中 $-1$ 出现两次。
            \textbf{理由：} 分解后的因式次数为 $1+1+2=4$，与原多项式次数一致。
    \end{description}

    \medskip
    \textbf{\textcolor{CExample}{例二（非首一：有理根的分母不一定是 $1$）}}
    \textbf{\textcolor{CExample}{题目：}}
    分解 $6x^3+7x^2-x-2$，求全部有理根。
    \begin{description}[style=nextline, leftmargin=2.8em, labelsep=0.5em, itemsep=0.46em]
        \item[\textbf{第一步（列候选值）：}]
            由 $p\mid(-2)$、$q\mid6$，可取 $p=\pm1,\pm2$，$q=1,2,3,6$，仅保留最简分数。先试 $x=1/2$。
            \textbf{理由：} 最高次项系数不是 $\pm1$，不能只试整数。
        \item[\textbf{第二步（验算根）：}]
            \[
               f\!\left(\frac12\right)
               =6\cdot\frac18+7\cdot\frac14-\frac12-2
               =\frac34+\frac74-\frac12-2=0.
            \]
            故 $2x-1$ 是因式。
        \item[\textbf{第三步（除法分解）：}]
            按方法四的逐列消元图，除以 $2x-1$ 的商为 $3x^2+5x+2$，乘回核验：
            \[
               (2x-1)(3x^2+5x+2)=6x^3+7x^2-x-2.
            \]
            继续分解 $3x^2+5x+2=(3x+2)(x+1)$。
        \item[\textbf{第四步（得出全部根）：}]
            \[
              \boxed{f(x)=(2x-1)(3x+2)(x+1)},
            \]
            所以 $x=\dfrac12,-\dfrac23,-1$。
            \textbf{理由：} 用 $qx-p$ 保留整数系数，比先写 $x-p/q$ 更方便。
    \end{description}

    \medskip
    \textbf{\textcolor{CExample}{例三（反向应用：排除一切有理根）}}
    \textbf{\textcolor{CExample}{题目：}}
    判断方程 $x^3-x-1=0$ 是否有有理根；它是否有实根？
    \begin{description}[style=nextline, leftmargin=2.8em, labelsep=0.5em, itemsep=0.45em]
        \item[\textbf{第一步（枚举候选）：}]
            $a_3=1$、$a_0=-1$，所以候选有理根只能是 $1$ 和 $-1$。
        \item[\textbf{第二步（逐个排除）：}]
            $f(1)=1-1-1=-1\ne0$；$f(-1)=-1+1-1=-1\ne0$。所以\textbf{没有有理根}。
        \item[\textbf{第三步（区分无有理根与无实根）：}]
            $f(1)=-1<0$，$f(2)=8-2-1=5>0$。因多项式函数连续，由零点存在定理，$(1,2)$ 内至少有一个实根。这个根不可能是有理数，因而是无理数。
            \textbf{理由：} 有理根定理只能排除有理数，不能排除无理实根。
    \end{description}

    \medskip
    \textbf{\textcolor{CExample}{例四（常数项为零：必须先提取公因式）}}
    \textbf{\textcolor{CExample}{题目：}}
    求 $2x^3-3x^2+x=0$ 的所有根。
    \begin{description}[style=nextline, leftmargin=2.8em, labelsep=0.5em, itemsep=0.45em]
        \item[\textbf{第一步（先取出 $x$）：}]
            \[
               2x^3-3x^2+x=x(2x^2-3x+1).
            \]
            因此 $x=0$ 已经是一个根。
        \item[\textbf{第二步（分解剩余部分）：}]
            \[
               2x^2-3x+1=(2x-1)(x-1).
            \]
            所以
            \[
               \boxed{x(2x-1)(x-1)=0\;\Longrightarrow\;x=0,\frac12,1.}
            \]
            \textbf{理由：} 对于 $a_0=0$ 的情形，不要遗漏零根，也不要把 $p=0$ 当成一个可用于“整除”的除数。
    \end{description}
\end{examplebox}
